% Blank lecture file. To add a lecture:
%   1. copy this file to tex/lecture_YYYY-MM-DD.tex and fill it in,
%   2. add these two lines at the end of the lecture list in main.tex:
%        \newlecture{Month D, YYYY}{Topic of the lecture}
%        \input{tex/lecture_YYYY-MM-DD}
% This file itself is not included in main.tex.

In this lecture, we \dots % one or two sentences summarizing the lecture

\section{First Topic}

Text of the notes goes here.

\begin{definition}[Name of the Definition]
    Statement of the definition.
\end{definition}

\begin{theorem}[Name of the Theorem]\label{thm:example}
    Statement of the theorem.
\end{theorem}

\begin{proof}
    Proof of the theorem.
\end{proof}

% A figure drawn with TikZ.
\begin{figure}[ht]
    \centering
    \begin{tikzpicture}
        \draw (0,0) rectangle (8,0.7);
        \foreach \x in {1,2,3,7} {\draw (\x,0) -- (\x,0.7);}
        \node at (5,0.35) {$\cdots$};
        \node[below] at (0.5,0) {\small $1$};
        \node[below] at (7.5,0) {\small $n$};
    \end{tikzpicture}
    \caption{Caption of the figure.}
    \label{fig:example}
\end{figure}

% An exercise given in class (dashed box in my notes), with space for the solution.
\begin{exercise}
    Statement of the exercise.
\end{exercise}
\solutionspace % TODO: replace with \begin{proof} ... \end{proof}
